Se compararon una CNN de referencia y YOLO26n-cls sobre una partición común de imágenes renales, con cinco semillas por arquitectura. La auditoría identificó 1.225 copias exactas adicionales y permitió conservar 10.531 imágenes con agrupamiento de similitud antes de la partición. YOLO26n-cls obtuvo el mayor F1 macro medio en el experimento, mientras que CNN presentó la menor latencia mediana media. El efecto medio YOLO menos CNN fue 0{,}0400 y el valor p bilateral exacto fue 0{,}1875.

La elección práctica debe considerar el F1 de cada clase, la dispersión entre corridas y el costo por imagen. La resolución estadística de cinco pares y la ausencia de identificadores verificables de paciente limitan la fuerza de las conclusiones. Como trabajo futuro se plantea recuperar procedencia y agrupación clínica, realizar validación externa y estudiar configuraciones de transferencia con presupuesto equivalente. El repositorio local conserva manifiesto, código, configuraciones y predicciones para reproducir la comparación.